% *** ก่อนส่ง: ลบข้อความตัวอย่าง (\ldots) ทั้งหมด แล้วเขียนเนื้อหาของตนเอง ***
% =============================================================
% บทที่ 5 - บทสรุป
% แก้ไขเนื้อหาในบทนี้ได้ที่ไฟล์ chapter5.tex
% =============================================================

% -----------------------------------------------------------
\section{สรุปผลการดำเนินโครงงาน}
% -----------------------------------------------------------

โครงงานนี้ได้พัฒนาระบบ\ldots{} ซึ่งมีวัตถุประสงค์เพื่อ\ldots{} ผลการดำเนินงานสรุปเทียบกับวัตถุประสงค์ได้ดังตารางที่ \ref{tab:obj_result}

\begin{table}[H]
\centering
\caption{สรุปผลเทียบกับวัตถุประสงค์}
\label{tab:obj_result}
\begin{tabular}{|c|l|l|c|}
\hline
\textbf{ข้อ} & \textbf{วัตถุประสงค์} & \textbf{ผลที่ได้} & \textbf{สถานะ} \\
\hline
1 & เพื่อศึกษาและวิเคราะห์\ldots & ศึกษาและวิเคราะห์\ldots{} แล้ว & บรรลุ \\ \hline
2 & เพื่อออกแบบและพัฒนา\ldots & พัฒนาระบบ\ldots{} เสร็จสมบูรณ์ & บรรลุ \\ \hline
3 & เพื่อทดสอบและประเมิน\ldots & ทดสอบด้วย UAT ผ่านทุกข้อ & บรรลุ \\
\hline
\end{tabular}
\end{table}

จากตารางที่ \ref{tab:obj_result} โครงงานนี้บรรลุวัตถุประสงค์ทั้ง\ldots{} ข้อ

% -----------------------------------------------------------
\section{ข้อจำกัดของระบบ}
% -----------------------------------------------------------

\subsection{ข้อจำกัดด้านการทำงาน}

ระบบยังไม่รองรับ\ldots{}

\subsection{ข้อจำกัดด้านฮาร์ดแวร์และซอฟต์แวร์}

ระบบต้องการ\ldots{} ในการทำงาน

% -----------------------------------------------------------
\section{ปัญหาอุปสรรค และแนวทางแก้ไข}
% -----------------------------------------------------------

ปัญหาที่พบระหว่างดำเนินงานและแนวทางแก้ไขแสดงดังตารางที่ \ref{tab:problems}

\begin{table}[H]
\centering
\caption{ปัญหาอุปสรรคและแนวทางแก้ไข}
\label{tab:problems}
\begin{tabular}{|c|l|l|}
\hline
\textbf{ลำดับ} & \textbf{ปัญหาที่พบ} & \textbf{แนวทางแก้ไข} \\
\hline
1 & \ldots & \ldots \\ \hline
2 & \ldots & \ldots \\ \hline
3 & \ldots & \ldots \\
\hline
\end{tabular}
\end{table}

% -----------------------------------------------------------
\section{ข้อเสนอแนะ ในการพัฒนาต่อไป}
% -----------------------------------------------------------

\subsection{การปรับปรุงฟังก์ชันการทำงาน}

ควรเพิ่มความสามารถในการ\ldots{}

\subsection{แนวทางการวิจัยต่อยอด}

สามารถนำไปต่อยอดในด้าน\ldots{}
